\documentclass[11pt]{article}
\usepackage[margin=1in]{geometry}
\usepackage{amsmath,amssymb,bm,booktabs,enumitem}
\usepackage{algorithm,algpseudocode}
\usepackage[hidelinks]{hyperref}

\newcommand{\R}{\mathbb{R}}
\newcommand{\C}{\mathbb{C}}
\newcommand{\vb}[1]{\mathbf{#1}}
\newcommand{\uv}[1]{\hat{\mathbf{#1}}}
\newcommand{\jj}{\mathrm{j}}
\newcommand{\Id}{\mathrm{Id}}
\renewcommand{\Re}{\operatorname{Re}}
\renewcommand{\Im}{\operatorname{Im}}

\title{Field Vectoring with a Four-Channel Coil Array\\[2pt]\large Lecture notes on \texttt{lib/VectorField}}
\author{ESP32\_PMW}
\date{2026-11-11}

\begin{document}
\maketitle

\begin{abstract}
We derive the algorithm in \texttt{VectorField::settingsFor}. Given a point $\vb r$, an
orientation $R$ and a field strength $b$, it returns each channel's carrier duty and phase so that
the field at $\vb r$ rotates in a circle of radius $b$ about the axis $R\uv z$. We then derive the
cost of tilting the axis, the effect of a twist about the axis, the effect of calibration errors,
and the maximum field the current limits allow for each axis. The longer write-up is
\texttt{controller/control/theory.md} \S12 and \S27.
\end{abstract}

\section{Problem statement}

The rotor's spin axis follows the rotation axis of the field at its position, and the rotor's
thrust points along its spin axis. To tilt the thrust we must set the field's rotation axis.

Per channel, the firmware sets a carrier duty, which sets the current amplitude, and a phase. All
four channels share one frequency. We want the duties and phases that make, at a point $\vb r$, a
circular field of strength $b$ rotating about a unit axis $\uv n$.

\paragraph{Notation.} Bold symbols are vectors, hats are unit vectors, and $\jj$ is the
imaginary unit ($\jj^2=-1$), which frees $i$ for currents. $\Id_3$ is the $3\times3$ identity.
$\vb r=(x,y,z)$ is a point in the array frame.

\section{Model of the array}

\paragraph{Geometry.} The four channels A, B, C, D sit at azimuths $0^\circ$, $90^\circ$,
$180^\circ$ and $270^\circ$ (measured 2026-08-29). Each channel is a $2\times2$ block of air-core
coils, centred $37$\,mm from the array axis, with $21$\,mm between coil centroids. So channel A's
coils sit at $(37\pm10.5,\ \pm10.5,\ 0)$\,mm. Every coil has radius $a=10.5$\,mm, $N=650$ turns and
its normal along $+\uv z$. The coil plane is $z=0$. The robot flies at $z$ between about 15 and
25\,mm, near the array axis.

\paragraph{Currents.} All four coil currents are sinusoids at one shared drive frequency $f$, so
they all repeat with period $1/f$. We call one such period a \emph{drive cycle}. To say where we
are inside a drive cycle, we define the \emph{drive angle}
\begin{equation}
  \theta(t) = 2\pi f t ,
\end{equation}
which advances by $2\pi$ every drive cycle.

Channel $k$'s current has its own amplitude and its own phase offset from $\theta$. We describe
it by a complex number $I_k$, the channel's \emph{phasor}, and write
\begin{equation}
  i_k(t) = \Re\!\big(I_k e^{\jj\theta}\big)
  = \Re\!\big(|I_k|\,e^{\jj(\theta + \arg(I_k))}\big)
  = |I_k|\cos\!\big(\theta + \arg(I_k)\big).
\end{equation}
Here $|I_k| = \sqrt{\Re(I_k)^2 + \Im(I_k)^2}$ is the magnitude of $I_k$, its length in the
complex plane, and it is the current amplitude in amps. $\arg(I_k)$ is the angle of $I_k$ in the
complex plane, measured from the positive real axis. The current peaks at $\theta = -\arg(I_k)$,
so $\arg(I_k)$ is how far channel $k$ runs ahead of $\theta$. The four phasors form
$\vb I = (I_A, I_B, I_C, I_D)\in\C^4$.

\paragraph{Field per amp.} We replace each coil by a point dipole. Let $\vb c$ be the coil's
centroid; we place the dipole at $\vb c$ and give it the coil's moment per amp,
$\vb m = N\pi a^2\,\uv z$. The robot sits about 35 to 55\,mm from the coil centroids, several coil
radii away, where the dipole term dominates the field of a loop. Section~7.2 compares this model
with the exact loop model. A channel's field per amp at $\vb r$ is the sum over its four coils,
\begin{equation}
  \vb b_k(\vb r) = \frac{\mu_0}{4\pi}\sum_{\vb c\in k}\left[
  \frac{3(\vb m\cdot\vb d)\,\vb d}{|\vb d|^5} - \frac{\vb m}{|\vb d|^3}\right],
  \qquad \vb d = \vb r - \vb c .
\end{equation}
We stack the four columns into the \emph{channel matrix}
\begin{equation}
  B(\vb r) = [\,\vb b_A\ \ \vb b_B\ \ \vb b_C\ \ \vb b_D\,]\in\R^{3\times4}\quad(\text{T/A}).
\end{equation}
Column $k$ is the field that one amp in channel $k$ makes at $\vb r$. Fields superpose, so
currents $\vb i$ make the field $B\vb i$. At $f\le 250$\,Hz the wavelength is kilometres and
nothing nearby conducts, so the field is quasi-static and $B$ does not depend on $f$. From here
on, $B$ is evaluated at the robot's position.

\section{Field over one drive cycle}

Substituting the currents into $\vb B = B\vb i$,
\begin{equation}
  \vb B(\theta) = \sum_k \vb b_k\, i_k
  = \Re\!\big(B\vb I\,e^{\jj\theta}\big)
  = B\,\Re(\vb I)\,\cos\theta - B\,\Im(\vb I)\,\sin\theta .
\end{equation}
In the last step we expand $e^{\jj\theta}=\cos\theta+\jj\sin\theta$ and use that $B$ is real. We
name the two coefficients:
\begin{equation}
  \boxed{\ \vb B(\theta) = \vb u\cos\theta + \vb v\sin\theta,\qquad
  \vb u = B\,\Re(\vb I),\qquad \vb v = -B\,\Im(\vb I)\ }
\end{equation}
$\vb u$ is the field at $\theta=0$ and $\vb v$ is the field at $\theta=\pi/2$. The tip of $\vb B$
traces an ellipse in the plane spanned by $\vb u$ and $\vb v$.

\paragraph{Rotation axis.} We differentiate, $d\vb B/d\theta = -\vb u\sin\theta + \vb v\cos\theta$,
and take the cross product:
\begin{equation}
  \vb B\times\frac{d\vb B}{d\theta}
  = \cos^2\theta\,(\vb u\times\vb v) - \sin^2\theta\,(\vb v\times\vb u)
  = \vb u\times\vb v .
\end{equation}
The result does not depend on $\theta$, so the field turns about the fixed direction
\begin{equation}
  \uv n = \frac{\vb u\times\vb v}{\lVert\vb u\times\vb v\rVert},
\end{equation}
by the right-hand rule. The sign of $\uv n$ gives the spin sense.

\paragraph{Circular field.} $|\vb B(\theta)|^2 = |\vb u|^2\cos^2\theta + |\vb v|^2\sin^2\theta
+ 2\,\vb u\cdot\vb v\,\sin\theta\cos\theta$. This equals $b^2$ for every $\theta$ if and only if
$|\vb u| = |\vb v| = b$ and $\vb u\cdot\vb v = 0$. Then the field has constant magnitude $b$ and
turns at the constant rate $2\pi f$.

\section{Commanded orientation}

A circular field of radius $b$ is fixed by $\vb u$ and $\vb v$: two perpendicular vectors of
length $b$. Their directions $\uv e_1 = \vb u/b$, $\uv e_2 = \vb v/b$ and
$\uv n = \uv e_1\times\uv e_2$ form a right-handed orthonormal frame, which is the columns of a
rotation matrix. So the caller supplies a rotation $R\in SO(3)$, or a unit quaternion that the
library converts to one, and we read its columns as
\begin{align}
  \uv e_1 &= R\uv x &&\text{field direction at }\theta=0,\nonumber\\
  \uv e_2 &= R\uv y &&\text{field direction at }\theta=\pi/2,\\
  \uv n &= R\uv z &&\text{rotation axis.}\nonumber
\end{align}
The target is $\vb u = b\,\uv e_1$, $\vb v = b\,\uv e_2$. Then
$\vb u\times\vb v = b^2\,\uv e_1\times\uv e_2 = b^2\uv n$, because $\det(R) = +1$. $R=\Id_3$ asks for a
field that starts along $+x$ and rotates about $+z$.

\paragraph{Twist about the axis.} $R$ has three angles and the axis $\uv n$ has two. We find what
the third does. We twist $R$ by an angle $\psi$ about its own axis, $R' = R\,R_z(\psi)$. The new
columns are
\begin{equation}
  \uv e_1' = \cos\psi\,\uv e_1 + \sin\psi\,\uv e_2,\qquad
  \uv e_2' = -\sin\psi\,\uv e_1 + \cos\psi\,\uv e_2,\qquad \uv n' = \uv n .
\end{equation}
The solution in Section~5 depends on $R$ only through $\uv e_1 - \jj\uv e_2$, which transforms as
\begin{equation}
  \uv e_1' - \jj\uv e_2'
  = (\cos\psi + \jj\sin\psi)\,\uv e_1 + (\sin\psi - \jj\cos\psi)\,\uv e_2
  = e^{\jj\psi}\,(\uv e_1 - \jj\uv e_2),
\end{equation}
since $\sin\psi - \jj\cos\psi = -\jj e^{\jj\psi}$. The solution is linear in this vector, so
$\vb I' = e^{\jj\psi}\vb I$. Each amplitude $|I_k|$ is unchanged and each phase advances by $\psi$.
In the field,
\begin{equation}
  \vb B'(\theta) = \Re\!\big(B\,e^{\jj\psi}\vb I\,e^{\jj\theta}\big) = \vb B(\theta+\psi),
\end{equation}
the same field shifted in time by $\psi/(2\pi f)$. The twist sets the time origin and nothing else.

\paragraph{Continuity of the twist.} The firmware applies a new phase at once, so a step
$\Delta\psi$ in the twist steps the field direction by $\Delta\psi$. A synchronous rotor that lags
the field by $\varphi$ makes torque $\propto\sin\varphi$, which peaks at $\varphi=\pi/2$, so it
slips once $\varphi+\Delta\psi>\pi/2$. The twist must therefore change continuously along a
trajectory. If the API took only $\uv n$, the library would have to choose $\uv e_1$ as a function
of $\uv n$, a tangent vector field on $S^2$, and by the hairy-ball theorem every such choice is
discontinuous somewhere. So the caller supplies the twist, and a quaternion slerp between
commands keeps it continuous.

\section{Inversion}

\subsection{Real linear system}
The field is nonlinear in amplitude and phase, because of the cosine, but linear in the real and
imaginary parts of $\vb I$. We set
\begin{equation}
  \vb x = \Re(\vb I),\qquad \vb y = -\Im(\vb I)\qquad(\vb x,\vb y\in\R^4),
\end{equation}
so that
\begin{equation}
  i_k(t) = \Re\big((x_k - \jj y_k)(\cos\theta + \jj\sin\theta)\big) = x_k\cos\theta + y_k\sin\theta .
\end{equation}
$\vb x$ is the cosine part of each channel's current and $\vb y$ the sine part. Section~3 gives
$\vb u = B\vb x$ and $\vb v = B\vb y$, and the target is
\begin{equation}
  \begin{bmatrix} B & 0\\ 0 & B\end{bmatrix}
  \begin{bmatrix}\vb x\\ \vb y\end{bmatrix}
  = \begin{bmatrix} b\,\uv e_1\\ b\,\uv e_2\end{bmatrix}.
\end{equation}
That is six equations in eight unknowns. The matrix is block diagonal, so it splits into two
independent problems, $B\vb x = b\uv e_1$ and $B\vb y = b\uv e_2$, each three equations in four
unknowns. Wherever $B$ has rank 3, each has a one-parameter family of solutions, and we need a
rule to pick one.

\subsection{Minimum-heat solution}
We pick the solution that dissipates the least power, because coil temperature is the binding
limit on the rig. The mean ohmic power in channel $k$ is
$R_s\langle i_k^2\rangle = \tfrac{R_s}{2}(x_k^2+y_k^2)$, with $R_s$ the series resistance. Summed
over channels it is $\tfrac{R_s}{2}\big(\lVert\vb x\rVert^2+\lVert\vb y\rVert^2\big)$. The two
blocks share no unknowns, so we minimise each norm separately. For one block with target $\vb t$
we solve
\begin{equation}
  \min_{\vb x}\ \tfrac12\lVert\vb x\rVert^2\quad\text{subject to}\quad B\vb x=\vb t .
\end{equation}
We form the Lagrangian $\mathcal L = \tfrac12\lVert\vb x\rVert^2 - \bm\lambda^\top(B\vb x-\vb t)$.
Setting $\partial\mathcal L/\partial\vb x = 0$ gives $\vb x = B^\top\bm\lambda$. Substituting this
into the constraint, we get $BB^\top\bm\lambda = \vb t$, so
\begin{equation}
  \vb x = P\,\vb t,\qquad P = B^\top(BB^\top)^{-1}\in\R^{4\times3},
\end{equation}
the pseudo-inverse of $B$. This is the global minimum. Any other solution is $\vb x+\bm\eta$ with
$B\bm\eta = 0$, and $\vb x^\top\bm\eta = \bm\lambda^\top B\bm\eta = 0$, so
$\lVert\vb x+\bm\eta\rVert^2 = \lVert\vb x\rVert^2 + \lVert\bm\eta\rVert^2 \ge \lVert\vb x\rVert^2$.

We apply $P$ to both blocks and recombine $\vb I = \vb x - \jj\vb y$:
\begin{equation}
  \boxed{\ \vb I = P(\vb u - \jj\vb v) = b\,P\,(\uv e_1 - \jj\uv e_2)\ }
\end{equation}
It is exact wherever $B$ has rank 3, because then $BP = BB^\top(BB^\top)^{-1} = \Id_3$. The
firmware inverts the $3\times3$ Gram matrix $G = BB^\top$ as $\operatorname{adj}(G)/\det(G)$.

\subsection{Heat as a function of the axis}
$P^\top P = G^{-1}BB^\top G^{-1} = G^{-1}$, so $\lVert P\vb t\rVert^2 = \vb t^\top G^{-1}\vb t$ and
\begin{equation}
  \lVert\vb x\rVert^2 + \lVert\vb y\rVert^2
  = b^2\big(\uv e_1^\top G^{-1}\uv e_1 + \uv e_2^\top G^{-1}\uv e_2\big)
  = b^2\big(\operatorname{tr}G^{-1} - \uv n^\top G^{-1}\uv n\big).
\end{equation}
The second step writes $\vb t^\top G^{-1}\vb t = \operatorname{tr}(G^{-1}\vb t\vb t^\top)$ and uses
$\uv e_1\uv e_1^\top + \uv e_2\uv e_2^\top = \Id_3 - \uv n\uv n^\top$. The heat depends on $\uv n$
and not on the twist. It is minimised by the $\uv n$ that maximises $\uv n^\top G^{-1}\uv n$,
the eigenvector of $G$ with the smallest eigenvalue $\sigma_3^2$.

\subsection{Solution on the array axis}
We write the singular value decomposition $B = \sum_j\sigma_j\,\vb a_j\vb w_j^\top$, with field
directions $\vb a_j\in\R^3$ and drive patterns $\vb w_j\in\R^4$. Then
$P = \sum_j\sigma_j^{-1}\vb w_j\vb a_j^\top$: field along $\vb a_j$ is made by the pattern
$\vb w_j$ at $1/\sigma_j$ amps per tesla.

On the array axis, $\vb r = (0,0,z)$, the four-fold symmetry fixes the form of $B$. Channel A's
coils lie at $+x$. A $+\uv z$ dipole at $(+x_c,0,0)$ gives $\vb d = (-x_c,0,z)$ and an
$x$-component $\propto 3z(-x_c)/|\vb d|^5<0$, so channel A's field at the axis points toward $-x$.
Rotating by $90^\circ$ for each channel,
\begin{equation}
  B = \begin{bmatrix} -\beta & 0 & \beta & 0\\ 0 & -\beta & 0 & \beta\\ \gamma&\gamma&\gamma&\gamma
  \end{bmatrix},
  \qquad
  G = BB^\top = \operatorname{diag}(2\beta^2,\ 2\beta^2,\ 4\gamma^2),
\end{equation}
where $\beta>0$ is one channel's inward field per amp and $\gamma$ its vertical field per amp. The
off-diagonal entries of $G$ vanish because $\gamma(-\beta+\beta) = 0$. $G$ is diagonal, so the
$\vb a_j$ are the coordinate axes, $\sigma_j = \sqrt{G_{jj}}$, and
$\vb w_j = B^\top\vb a_j/\sigma_j$:
\begin{center}
\begin{tabular}{lll}\toprule
$\vb a_j$ & $\sigma_j$ & $\vb w_j$\\\midrule
$\uv x$ & $\sqrt2\,\beta$ & $(-1,0,1,0)/\sqrt2$, A against C\\
$\uv y$ & $\sqrt2\,\beta$ & $(0,-1,0,1)/\sqrt2$, B against D\\
$\uv z$ & $2\gamma$ & $(1,1,1,1)/2$, the common mode\\
none & $0$ & $(1,-1,1,-1)/2$, the null mode\\\bottomrule
\end{tabular}
\end{center}
The null mode satisfies $B\vb w = 0$: it makes no field at $\vb r$ and the minimum-norm solution
contains none of it. At $z = 20.1$\,mm, $\sigma = (2.42,\ 2.42,\ 0.36)$\,mT/A, so vertical field
costs $2.42/0.36 = 6.7$ times more current per tesla than horizontal field.

The rows of $P = B^\top G^{-1}$ are
\begin{equation}
  \vb p_A = \Big(-\tfrac{1}{2\beta},\,0,\,\tfrac{1}{4\gamma}\Big),\quad
  \vb p_B = \Big(0,\,-\tfrac{1}{2\beta},\,\tfrac{1}{4\gamma}\Big),\quad
  \vb p_C = \Big(\tfrac{1}{2\beta},\,0,\,\tfrac{1}{4\gamma}\Big),\quad
  \vb p_D = \Big(0,\,\tfrac{1}{2\beta},\,\tfrac{1}{4\gamma}\Big),
\end{equation}
and $I_k = b\,\vb p_k\cdot(\uv e_1 - \jj\uv e_2)$.

\paragraph{Vertical axis.} $R = \Id_3$ gives $\uv e_1 - \jj\uv e_2 = (1,-\jj,0)$, so
\begin{equation}
  \vb I = \frac{b}{2\beta}\,(-1,\ \jj,\ 1,\ -\jj).
\end{equation}
The four phasors have equal magnitude and angles $180^\circ$, $90^\circ$, $0^\circ$ and $-90^\circ$:
each channel lags the previous one by $90^\circ$.

\paragraph{Tilted axis.} We tilt by $\alpha$ toward $+x$, taking $R$ as a rotation about $\uv y$:
$\uv e_1 = (\cos\alpha,0,-\sin\alpha)$, $\uv e_2 = \uv y$, $\uv n = (\sin\alpha,0,\cos\alpha)$. Then
\begin{equation}
  I_k = b\big(p_{k,x}\cos\alpha - \jj\,p_{k,y} - p_{k,z}\sin\alpha\big),
  \qquad
  \vb I = \frac{b}{2\beta}\,(-\cos\alpha,\ \jj,\ \cos\alpha,\ -\jj)
  \;-\;\frac{b\sin\alpha}{4\gamma}\,(1,1,1,1).
\end{equation}
The tilt puts a vertical component $-b\sin\alpha$ into $\vb u$, and only the common mode makes
vertical field, so every channel gains the same real term. Its size relative to the A--C drive is
\begin{equation}
  \frac{\sin\alpha/(4\gamma)}{\cos\alpha/(2\beta)} = \frac{\beta}{2\gamma}\tan\alpha
  \approx 4.8\tan\alpha\qquad(z=20.1\text{ mm}).
\end{equation}
At $\alpha = 10^\circ$ this ratio is $0.84$. From Section~5.3, with
$\operatorname{tr}G^{-1} = 1/\beta^2 + 1/(4\gamma^2)$, a vertical axis costs
$\operatorname{tr}G^{-1} - 1/(4\gamma^2) = 1/\beta^2$ and a horizontal axis costs
$\operatorname{tr}G^{-1} - 1/(2\beta^2) = 1/(2\beta^2) + 1/(4\gamma^2)$ per $b^2$. Their ratio is
$\tfrac12 + (\beta/2\gamma)^2 \approx 23$ at $20.1$\,mm.

\section{From currents to drive settings}

\paragraph{Amplitude.} Each H-bridge switches the coil between $+V$ and $-V$ at the drive
frequency, and a 20\,kHz carrier gates it with duty $D$. The RLC filters out the carrier, so the
coil sees a square wave of height $DV$. A square wave of height $A$ has fundamental amplitude
$\tfrac4\pi A$, and the RLC also attenuates the higher harmonics, so
\begin{equation}
  |I| = \frac{4DV}{\pi|Z(f)|},\qquad
  |Z(f)| = \sqrt{R_s^2 + \Big(\omega L_s - \frac{1}{\omega C_s}\Big)^2},\quad \omega=2\pi f
  \qquad\Longrightarrow\qquad
  D_k = \frac{\pi|Z(f)|\,|I_k|}{4V g_k},
\end{equation}
where $g_k$ is a per-channel gain trim, 1 until fitted. With $V = 12$\,V, $R_s = 6.9\,\Omega$,
$L_s = 1.4$\,mH and $C_s = 800\,\mu$F, the nominal resonance is $f_0\approx150$\,Hz; it has not
been measured on the current capacitor bank. At $f_0$, $|Z| = R_s$ and $D=100\%$ gives
$|I| = 2.2$\,A, so near resonance the $2$\,A continuous limit binds before the duty limit.

\paragraph{Saturation.} If any channel exceeds $100\%$ duty or $I_{cont}=2$\,A, the library scales
the whole set by
\begin{equation}
  s = \min\!\Big(1,\ \frac{100\%}{\max_k D_k},\ \frac{I_{cont}}{\max_k |I_k|}\Big).
\end{equation}
$\vb I\to s\vb I$ gives $\vb u\to s\vb u$ and $\vb v\to s\vb v$, so $\uv n$ and circularity are
unchanged and the field becomes $sb$. The library reports $s$ as \texttt{fieldFraction}.
Section~7.3 gives the policy flight needs instead.

\paragraph{Phase.} The firmware delays channel $k$ by $\varphi_k/360^\circ$ of a period, so its
phase setting is a lag, of opposite sign to $\arg(I_k)$. With a per-channel offset $o_k$,
\begin{equation}
  \varphi_k = o_k - \arg(I_k) \pmod{360^\circ}.
\end{equation}
The offsets absorb winding direction, wiring, and any phase the drive electronics add. The phase
pattern predicted from the azimuths alone does not spin the rotor, and the empirical
\texttt{SPIN\_PHASES} does, so we fit $o_k$ in the constructor: $R=\Id_3$ at $(0,0,15)$\,mm must
send exactly \texttt{SPIN\_PHASES}. The fit is $o = (90^\circ, 180^\circ, 180^\circ, 270^\circ)$.
A reversed winding appears as a $180^\circ$ offset. We cannot yet explain the $90^\circ$ parts, so
we treat the tilt direction as unverified until we run a step test on the rig.

\paragraph{Calibration errors.} We suppose the real currents are $\kappa_k I_k$ for some complex
$\kappa_k$, an error in impedance, gain or phase. If all four share one $\kappa$, the field is
\begin{equation}
  \Re\!\big(B\,\kappa\vb I\,e^{\jj\theta}\big) = |\kappa|\,\Re\!\big(B\vb I\,e^{\jj(\theta+\arg(\kappa))}\big)
  = |\kappa|\,\vb B(\theta+\arg(\kappa)),
\end{equation}
the commanded field scaled by $|\kappa|$ and shifted in time, with the same axis. If the
$\kappa_k$ differ, the field is $\Re(B\operatorname{diag}(\kappa)\vb I e^{\jj\theta})$, which in
general has a different axis. So the RLC values set the field magnitude only, and the axis depends
only on the channel-to-channel trims $g_k$ and $o_k$.

\section{Algorithm}

Algorithm~\ref{alg:settings} is \texttt{VectorField::settingsFor}. It never refuses with an
error. Every input gets a defined output, and the outputs that make no field are all
$D_k=0$. Channels are indexed $k=1,\dots,4$ (A, B, C, D), and every $\max_k$ is over those four.

\begin{algorithm}[H]
\caption{\texttt{VectorField::settingsFor}$(\vb r, R, b, f)$}\label{alg:settings}
\begin{algorithmic}[1]
\Require point $\vb r\in\R^3$ (m, array frame); orientation $R\in\R^{3\times3}$ or quaternion
  $q\in\R^4$; field radius $b\in\R$ (T, signed); drive frequency $f$ (Hz)
\Require per-channel gain trim $g_k$ (1 until fitted) and phase offset $o_k$ (deg, fitted to
  \texttt{SPIN\_PHASES}, Section~6)
\Ensure carrier duty $D_k$ (\%), phase lag $\varphi_k$ (deg), field fraction $s\in[0,1]$
\If{a quaternion was given}
  \If{$\lVert q\rVert = 0$}
    \State \Return $D_k = 0$, $s = 0$ \Comment{no orientation: the caller keeps its previous drive}
  \EndIf
  \State $R \gets R(q/\lVert q\rVert)$ \Comment{every nonzero $q$ gives a rotation}
\ElsIf{$RR^\top \ne \Id_3$ or $\det(R) \ne +1$ (each to $10^{-3}$)}
  \State \Return $D_k = 0$, $s = 0$ \Comment{a stretch or a mirror: the caller keeps its previous drive}
\EndIf
\If{$b$ is not finite}
  \State \Return $D_k = 0$, $s = 0$ \Comment{NaN or $\infty$ from the wire}
\EndIf
\State $\uv e_1 \gets R\uv x$, \quad $\uv e_2 \gets R\uv y$ \Comment{Section~4}
\State $B \gets [\,\vb b_A\ \vb b_B\ \vb b_C\ \vb b_D\,]$ at $\vb r$ \Comment{dipole sum, Section~2; T/A}
\State $G \gets BB^\top$ \Comment{$3\times3$ Gram matrix}
\If{$|\det(G)| \le 10^{-40}\ (\mathrm{T/A})^6$}
  \State \Return $D_k = 0$, $s = 0$ \Comment{$B$ has rank $<3$ at $\vb r$}
\EndIf
\State $G^{-1} \gets \operatorname{adj}(G)/\det(G)$
\State $\vb x \gets B^\top G^{-1}(b\,\uv e_1)$, \quad $\vb y \gets B^\top G^{-1}(b\,\uv e_2)$
  \Comment{minimum heat, Section~5.2}
\State $I_k \gets x_k - \jj y_k$, \quad $|I_k| \gets \sqrt{\Re(I_k)^2 + \Im(I_k)^2}$ \Comment{A}
\If{$f \le 0$}
  \State \Return $D_k = 0$, $s = 0$ \Comment{$C_s$ blocks DC: $|Z(f)|\to\infty$}
\EndIf
\State $\omega \gets 2\pi f$, \quad
  $|Z(f)| \gets \sqrt{R_s^2 + (\omega L_s - 1/(\omega C_s))^2}$ \Comment{$\Omega$}
\State $D_k \gets 100\%\cdot\pi\,|Z(f)|\,|I_k|\,/\,(4Vg_k)$ \Comment{Section~6}
\If{$\max_k |I_k| > 0$}
  \State $s \gets \min\!\big(1,\ 100\%/\max_k D_k,\ I_{cont}/\max_k |I_k|\big)$
\Else
  \State $s \gets 1$ \Comment{$b = 0$: every $D_k = 0$, so the coils are off}
\EndIf
\State $D_k \gets s\,D_k$
\State $\arg(I_k) \gets \operatorname{atan2}\!\big(\Im(I_k),\ \Re(I_k)\big)$
  \Comment{angle of $I_k$ from the positive real axis}
\State $\varphi_k \gets \big(o_k - \arg(I_k)\big) \bmod 360^\circ$, \ in $[0^\circ, 360^\circ)$
\State \Return $D_k$, $\varphi_k$ for $k = 1,\dots,4$, and $s$
\end{algorithmic}
\end{algorithm}

\paragraph{The safe outputs.} Two kinds of input make no field.
\begin{itemize}[nosep]
  \item An orientation that cannot be used ($q=0$, or $R\notin SO(3)$), a non-finite $b$, or a
        singular $B$ returns $s=0$. The firmware treats $s=0$ as ``do nothing'' and keeps the
        previous drive.
  \item $f\le0$ also returns $s=0$. This is the limit of the formula, because the series capacitor
        blocks DC, so $|Z|\to\infty$, $D_k\to\infty$ and $s\to0$. \texttt{freq=} already rejects
        $f\le0$, so this case cannot happen in FLIGHT.
\end{itemize}
The field radius $b$ is signed. The solve is linear in $b$, so $-b$ gives $\vb I\to-\vb I$, which
leaves each $|I_k|$ the same and adds $180^\circ$ to each $\arg(I_k)$. That is the field
$\vb B(\theta+\pi)$, the same as a twist $R\,R_z(\pi)$. So a sign flip in flight is a $180^\circ$
phase step under a locked rotor, which Section~4 says the rotor will not survive. $b=0$ returns
$D_k=0$ with $s=1$, and the firmware applies it, so the coils go off.
The FLIGHT-only command
\begin{center}
\texttt{vec=<x\_mm>:<y\_mm>:<z\_mm>:<qw>:<qx>:<qy>:<qz>:<b\_mT>}
\end{center}
calls this. The firmware stores the target and re-solves on every \texttt{freq=}, because the
duties depend on $|Z(f)|$ and the axis does not.

\section{Reachable workspace}

\subsection{Maximum field for a given axis}
From $I_k = b\,\vb p_k\cdot(\uv e_1 - \jj\uv e_2)$,
\begin{equation}
  |I_k|^2 = b^2\big[(\vb p_k\cdot\uv e_1)^2 + (\vb p_k\cdot\uv e_2)^2\big]
  = b^2\,\vb p_k^\top(\Id_3-\uv n\uv n^\top)\,\vb p_k
  = b^2\,|\vb p_k\times\uv n|^2 .
\end{equation}
The middle step uses $\uv e_1\uv e_1^\top + \uv e_2\uv e_2^\top = \Id_3 - \uv n\uv n^\top$; the
last uses $|\vb p|^2 - (\vb p\cdot\uv n)^2 = |\vb p\times\uv n|^2$. The twist does not appear.
Requiring $|I_k|\le I_{lim}$ on every channel gives
\begin{equation}
  \boxed{\ b_{max}(\vb r,\uv n) = \frac{I_{lim}}{\max_k|\vb p_k\times\uv n|}\ }
\end{equation}
where $I_{lim}$ is the smaller of $I_{cont}$ and the current at $100\%$ duty. Let $\delta_k$ be
the angle between $\uv n$ and $\vb p_k$, so $|\vb p_k\times\uv n| = |\vb p_k|\sin\delta_k$. For a
required $b$, channel $k$ allows the axes with $\sin\delta_k\le I_{lim}/(b|\vb p_k|)$, a double
cone about $\pm\vb p_k$. The reachable axes are the intersection of the four double cones.

On the array axis, a vertical axis gives $|\vb p_k\times\uv z| = 1/(2\beta)$ for every $k$, so
$b_{max} = 2\beta I_{lim}$. For the horizontal axis $\uv n=\uv x$,
$|\vb p_k\times\uv x|^2 = |\vb p_k|^2 - p_{k,x}^2$, which is largest for B and D:
\begin{equation}
  b_{max}(\uv x) = I_{lim}\Big/\sqrt{\frac{1}{4\beta^2} + \frac{1}{16\gamma^2}} .
\end{equation}
At $20.1$\,mm these give $6.8$ and $1.4$\,mT. Table~1 evaluates $b_{max}$ over height.

\begin{table}[h]\centering
\begin{tabular}{rccccc}\toprule
$z$ (mm) & vertical $b_{max}$ & horizontal $b_{max}$ & \multicolumn{3}{c}{largest tilt reachable at $b=$}\\
 & (mT) & (mT) & 2 mT & 3 mT & 4.5 mT\\\midrule
10 & 6.5 & 5.7 & $90^\circ$ & $90^\circ$ & $90^\circ$\\
15 & 7.3 & 4.5 & $90^\circ$ & $90^\circ$ & $90^\circ$\\
18 & 7.2 & 2.8 & $90^\circ$ & $44.5^\circ$ & $15^\circ$\\
20 & 6.9 & 1.5 & $34^\circ$ & $16.5^\circ$ & $6.5^\circ$\\
25 & 5.9 & 1.2 & $24^\circ$ & $11^\circ$ & $3.5^\circ$\\
30 & 4.8 & 2.3 & $90^\circ$ & $21^\circ$ & $2^\circ$\\\bottomrule
\end{tabular}
\caption{On the array axis, $150$\,Hz, dipole model, $I_{lim}=2$\,A. ``Largest tilt'' is the
half-angle of the largest cone about $\uv z$ whose every axis reaches $b$.}
\end{table}

The horizontal $b_{max}$ has a minimum between 20 and 25\,mm, set by the $1/(16\gamma^2)$ term.
Section~7.2 derives where $\gamma$ vanishes.

\subsection{The singular shell}
On the array axis $\sigma_3 = 2\gamma$. A $+\uv z$ dipole at lateral distance $\rho$ makes, at
height $z$ on the array axis, with $|\vb d|^2 = \rho^2 + z^2$ and $\vb m\cdot\vb d = mz$,
\begin{equation}
  B_z \propto \frac{3z^2}{|\vb d|^5} - \frac{1}{|\vb d|^3}
  = \frac{3z^2 - (\rho^2+z^2)}{(\rho^2+z^2)^{5/2}} = \frac{2z^2-\rho^2}{(\rho^2+z^2)^{5/2}},
\end{equation}
which is negative for $z<\rho/\sqrt2$ and positive above. Each channel's inner coils sit at
$\rho = \sqrt{26.5^2+10.5^2} = 28.5$\,mm and change sign at $20.2$\,mm. Its outer coils sit at
$\rho = \sqrt{47.5^2+10.5^2} = 48.6$\,mm and change sign at $34.4$\,mm. Between those heights the
two pairs have opposite sign, and their sum $\gamma$ is zero at $22.4$\,mm. Off the array axis we
find the zero of $\sigma_3$ numerically:

\begin{table}[h]\centering
\begin{tabular}{lccccc}\toprule
$x$ (mm) & 0 & 6 & 10 & 15 & 20\\\midrule
$z$ where $\sigma_3=0$, dipole model (mm) & 22.4 & 22.0 & 21.2 & 19.4 & 16.7\\
$z$ where $\sigma_3=0$, exact loop model (mm) & 20.8 & 20.3 & 19.4 & 17.5 & 15.0\\\bottomrule
\end{tabular}
\end{table}

Where $\sigma_3 = 0$, $B$ has rank 2 and every field it makes lies in the plane of $\vb a_1$ and
$\vb a_2$. A circular field must then have $\vb u$ and $\vb v$ in that plane, so
$\uv n = \pm\vb a_3$, which is $\pm\uv z$ on the array axis. This holds for any current. Near the
shell the solve succeeds, but $\vb p_k$ contains $1/(4\gamma)$, so the common mode grows without
bound and $s$ shrinks the field toward zero.

The two models place the shell $1.6$\,mm apart, and the real coils are $22$\,mm-long solenoids,
which neither model describes. We have to measure the height on the rig, by finding where a
commanded tilt stops tilting.

\subsection{Consequences for planning}
The rotor stays phase-locked only if
$b\ge b_{lock}(f) = 1.0\,\text{mT}\,(f/110\,\text{Hz})^2/0.8$ (theory.md \S12.3: drag torque grows
as $f^2$, with a 0.8 safety factor on the lock angle). That is $1.05$\,mT at $101$\,Hz and
$4.56$\,mT at $210$\,Hz. A pose $(\vb r,\uv n)$ is feasible at frequency $f$ when
$b_{lock}(f)\le b_{max}(\vb r,\uv n)$. From this:
\begin{enumerate}[nosep]
  \item We cross the shell with the axis vertical. A climb through it must bring the tilt in
        before the band and restore it after, so the planner must look ahead along the
        trajectory.
  \item In flight, we reduce tilt before amplitude. Scaling $b$ below $b_{lock}$ makes the rotor
        slip. We keep $b\ge b_{lock}$ and move $\uv n$ along the great circle toward $\uv z$ until
        $b_{max}\ge b$; $b_{max}$ is closed-form, so this is a one-dimensional bisection. The
        library's uniform scaling does the opposite.
  \item We limit the axis slew rate. The rotor's axis follows commands at about $0.5$\,Hz
        (theory.md \S12.6).
\end{enumerate}

\paragraph{Underactuation.} Away from the shell there are eight real unknowns and six field
components, so the field at a point is fully actuated. With the field fixed, the remaining two
unknowns are the complex amplitude of the null mode, which cannot set all three components of the
gradient force. Lateral position must therefore be controlled through thrust tilt, and the
Earnshaw anti-trap (theory.md \S12.6) is a disturbance set by the field choice.

\section{Verification and open items}
\begin{itemize}[nosep]
  \item \texttt{lib/VectorField/check/self\_check.cpp} checks:
        the geometry against the GUI ($4.58$ and $4.29$\,mT);
        $\sigma=(2.42,2.42,0.36)$\,mT/A at $20.1$\,mm and the null mode;
        200 random orientations reached to $10^{-3}$\,deg with a circular field;
        $R=\Id_3$ sending \texttt{SPIN\_PHASES};
        a twist changing only the phases;
        saturation keeping the axis;
        safe outputs for bad input (signed $b$, $b=0$, $f=0$, $q=0$, non-rotations).
  \item The C++ matched the original Python solver to float precision on 40 cases.
  \item Open, all needing the rig: the tilt direction (the offsets $o_k$), the shell height, $f_0$,
        and the gains $g_k$.
\end{itemize}

\end{document}
